%=============================================================================
% Phase 2 IEEE report -- Set-Associative Write-Back Cache
% Compile: pdflatex Phase2_Report.tex  (needs IEEEtran; Overleaf -> pdfLaTeX)
% Fill in every [Name / Roll No] before submitting.
%=============================================================================
\documentclass[conference]{IEEEtran}
\IEEEoverridecommandlockouts
\usepackage{cite,amsmath,booktabs,array,url}

\begin{document}

\title{From Direct-Mapped to Set-Associative: A Write-Back\\
Cache with Design-Space Exploration (Phase 2)}

\author{
\IEEEauthorblockN{Divy Kumar Jain}\IEEEauthorblockA{[Roll No]}\and
\IEEEauthorblockN{Praveen Kumar Nitharwal}\IEEEauthorblockA{[Roll No]}\and
\IEEEauthorblockN{Piyush Raj}\IEEEauthorblockA{[Roll No]}\and
\IEEEauthorblockN{Utkarsh Jain}\IEEEauthorblockA{[Roll No]}
}
\maketitle

\begin{abstract}
Phase 1 delivered a verified direct-mapped, read-only cache. Phase 2 completes
the project by extending it into a \emph{2-way set-associative, write-back}
cache in Verilog, and by exploring the cache design space. The new design adds
LRU replacement, a write path with a dirty bit and write-allocate, and a
writeback-on-eviction sequence, organised as 32 sets of 2 ways (tag 23, index
5, offset 4). A read/write trace-driven testbench and a parameterized Python
model (which agree exactly on hit ratios and writeback counts) measure the
effect of associativity, block size, cache size, and write policy. The central
result is that associativity eliminates conflict misses: on a conflict-heavy
trace the direct-mapped cache achieves 0\% hits while the 2-way cache achieves
99.22\%. A Yosys synthesis pass reports the cache's cell/area cost.
\end{abstract}

\begin{IEEEkeywords}
set-associative cache, LRU replacement, write-back, write-allocate, AMAT,
conflict miss, Verilog, logic synthesis.
\end{IEEEkeywords}

\section{Introduction}
A direct-mapped cache maps each block to exactly one line, so two frequently
used blocks that share a line evict one another repeatedly (conflict misses).
Set-associative caches give each set several ways, allowing such blocks to
coexist. Phase 2 implements a 2-way set-associative cache, adds write support
with a write-back policy, and quantifies the architectural trade-offs through a
parameter sweep and synthesis.

\section{Proposed Solution}
\subsection{Organisation and Mapping}
The cache keeps the Phase 1 capacity (1~KB, 16~B blocks) but reorganises it as
32 sets of 2 ways. The tag is compared against both ways of the indexed set in
parallel; a hit in either way returns the word in one cycle.

\subsection{LRU Replacement}
With two ways, one bit per set records the least-recently-used way. On any hit
or fill of way $w$, the other way becomes the victim ($lru \leftarrow \lnot w$);
on a miss the victim is the current LRU way.

\subsection{Write-Back Policy}
Write hits update the cached word and set a dirty bit (no memory traffic).
Write misses use write-allocate: the block is fetched, the store word merged,
and the line marked dirty. Evicting a dirty line writes it back to memory
before the new block is fetched. The controller FSM has four states: IDLE
(1-cycle hit), WB (write back a dirty victim), ALLOC (fetch the new block), and
FIN (answer the CPU).

\subsection{Verification}
A read/write testbench drives the RTL and counts read/write hits, misses, and
writebacks. An independent Python model reproduces the same associativity, LRU,
and write-back behaviour; exact agreement on hit ratio and writeback count is
the correctness criterion.

\section{Results}
\subsection{RTL vs.\ Reference Model}
Both produce identical hit ratios and writeback counts (Table~\ref{tab:rtl}).

\begin{table}[htbp]
\caption{RTL and model agreement (256 accesses, 2-way write-back)}
\label{tab:rtl}
\centering
\begin{tabular}{@{}lrrr@{}}
\toprule
\textbf{Trace} & \textbf{Hit \%} & \textbf{Writebacks} & \textbf{AMAT} \\
\midrule
conflict   & 99.22 & 0  & 1.078 \\
repeated   & 96.88 & 0  & 1.312 \\
sequential & 75.00 & 0  & 3.500 \\
random     & 17.97 & 40 & 9.203 \\
\bottomrule
\end{tabular}
\end{table}

\subsection{Design-Space Sweep}
Table~\ref{tab:sweep} compares configurations on the conflict and random
traces.

\begin{table}[htbp]
\caption{Hit ratio across configurations}
\label{tab:sweep}
\centering
\begin{tabular}{@{}lrr@{}}
\toprule
\textbf{Configuration} & \textbf{conflict} & \textbf{random} \\
\midrule
Direct-mapped 1\,KB/16\,B & 0.00  & 14.45 \\
2-way 1\,KB/16\,B         & 99.22 & 17.97 \\
4-way 1\,KB/16\,B         & 99.22 & 19.14 \\
Direct-mapped 1\,KB/32\,B & 0.00  & 20.70 \\
Direct-mapped 2\,KB/16\,B & 99.22 & 28.12 \\
\bottomrule
\end{tabular}
\end{table}

Three effects stand out: (i) associativity removes the conflict misses that
cripple the direct-mapped cache (0\% $\rightarrow$ 99.22\%); (ii) larger blocks
raise spatial-locality hit rates (sequential 75\% $\rightarrow$ 87.5\% at 32~B,
not shown); and (iii) more capacity reduces capacity misses on the random
trace. Write-back incurs writebacks only on dirty evictions (40 on the random
trace), far fewer memory writes than write-through would generate.

\subsection{Synthesis}
Synthesizing \texttt{sa\_cache} with Yosys~\cite{yosys} (\texttt{make synth})
maps the design to generic gates (Table~\ref{tab:synth}). The $\approx$9{,}800
flip-flops correspond almost entirely to the storage: two ways of 32 lines,
each 128-bit data $+$ 23-bit tag $+$ valid $+$ dirty ($\approx$153 bits), plus
the per-set LRU bits and the FSM state. The large mux count reflects
way-selection, word-selection, and writeback datapaths.

\begin{table}[htbp]
\caption{Yosys synthesis of \texttt{sa\_cache} (generic gates)}
\label{tab:synth}
\centering
\begin{tabular}{@{}lr@{}}
\toprule
\textbf{Resource} & \textbf{Count} \\
\midrule
Total cells        & 25{,}233 \\
Flip-flops (DFFE/DFF/SDFFE) & 9{,}828 \\
Multiplexers       & 8{,}786 \\
Wires              & 14{,}660 \\
\bottomrule
\end{tabular}
\end{table}

\section{Conclusion}
Phase 2 completes the project: a working 2-way set-associative write-back cache,
verified against a reference model, with a parameter sweep and synthesis that
quantify the cost and benefit of associativity, block size, capacity, and write
policy. The results confirm, on real traces, why modern caches are
set-associative.

\begin{thebibliography}{1}
\bibitem{hennessy} J.~L. Hennessy and D.~A. Patterson, \emph{Computer
Architecture: A Quantitative Approach}, 6th~ed. Morgan Kaufmann, 2019.
\bibitem{patterson} D.~A. Patterson and J.~L. Hennessy, \emph{Computer
Organization and Design}, 5th~ed. Morgan Kaufmann, 2014.
\bibitem{smith} A.~J. Smith, ``Cache memories,'' \emph{ACM Computing Surveys},
vol.~14, no.~3, pp.~473--530, 1982.
\bibitem{yosys} C.~Wolf, ``Yosys Open SYnthesis Suite,''
\url{https://yosyshq.net/yosys/}.
\end{thebibliography}

\end{document}
